% Generated draft: requires mathematical and source review.
\documentclass[11pt]{article}
\usepackage[margin=25mm]{geometry}
\usepackage{fontspec}
\setmainfont{Latin Modern Roman}
\usepackage{amsmath,amssymb,mathtools,unicode-math}
\setmathfont{Latin Modern Math}
\usepackage{enumitem,needspace,xurl,hyperref}
\setlength{\parindent}{0pt}
\setlength{\parskip}{5pt}
\setlength{\emergencystretch}{3em}
\newcommand{\N}{\mathbb{N}}
\newcommand{\Z}{\mathbb{Z}}
\newcommand{\Q}{\mathbb{Q}}
\newcommand{\R}{\mathbb{R}}
\newcommand{\C}{\mathbb{C}}
\newcommand{\F}{\mathbb{F}}
\newcommand{\A}{\mathbb{A}}
\newcommand{\PP}{\mathbb{P}}
\newcommand{\E}{\mathbb{E}}
\newcommand{\eps}{\varepsilon}
\newcommand{\dd}{\,\mathrm{d}}
\begin{document}
\section*{UnsolvedMath ID: 30000551}
\textbf{Source problem code:} catalogue entry
\subsection{Definitions and mathematical statement}
Under the smooth Riemannian interpretation, let $g$ be a complete Riemannian metric on $\R^2$, and let $K_g$ be its Gaussian curvature. Flatness means $K_g\equiv0$; for a complete simply connected Riemannian surface this is equivalent to being globally isometric to the Euclidean plane.

One usual formulation of the unique-parallel-geodesic property is the Euclidean parallel axiom: if $\mathcal L_g$ is the intended collection of unparametrized complete geodesic images, then
\[
\forall\ell\in\mathcal L_g\ \forall x\in\R^2\setminus\ell,
\quad\exists!\,m\in\mathcal L_g:
\quad x\in m\ \text{ and }\ m\cap\ell=\varnothing.
\]
Uniqueness here concerns images, so changing the origin or orientation of a geodesic parametrization does not create another parallel.

The source asks whether completeness and its unique-parallel-geodesic property imply $K_g\equiv0$. The displayed parallel axiom is an explicit conventional interpretation, not a recovered source definition. The extract does not specify whether ``parallel'' means disjointness, bounded distance, or an asymptotic relation, nor whether the geodesics are all complete geodesics or only globally minimizing lines. It also does not explicitly say that $g$ is smooth Riemannian rather than a more general metric. These distinctions must be resolved before the implication is a uniquely determined theorem; no additional no-conjugate-points or curvature-bound hypothesis is inferred.
\subsection{Short English statement}
Must a complete metric on the plane be flat if each geodesic has exactly one parallel through each point outside it? The source leaves the metric category and the meaning of parallel and geodesic unspecified; the usual smooth disjoint-geodesic interpretation is stated explicitly.
\subsection{Sources}
\begin{itemize}
\item Supplied report: \emph{Geometric Group Theory, Hyperbolic Dynamics and Symplectic Geometry} (2007), Oberwolfach Reports, DOI 10.4171/owr/2006/33, pp. 1991--2058; proposed year 2006. \url{https://doi.org/10.4171/owr/2006/33}.
\item Local UnsolvedMath registry, permanent ID 30000551, catalogue entry; created 25 August 2026, updated 16 September 2026. The corrected\_verified note describes editorial cleanup. The supplied literature triage of 21 August 2026 labels the fragment open conservatively because surrounding hypotheses on $g$ are missing, rather than reporting a verified fully specified open theorem. No further external reference is supplied; the dated assessment is retained without independent verification.
\end{itemize}
\end{document}
